\subsection{نتایج تجربی}
ارزیابی ماژول \lr{qcqa\_attention} بر روی بنچ‌مارک استاندارد \lr{GLUE SST-2} با طول دنباله ۵۱۲ بر روی پردازنده گرافیکی در شرایط کاملاً کنترل‌شده انجام گرفته است. داده‌های رسمی استخراج‌شده از \lr{master\_verified\_results.json} در جدول \ref{tab:benchmark_qcqa_attention} درج شده است:

\begin{table}[H]
\centering
\small
\caption{کارنامه تجربی و ارزیابی بنچ‌مارک سخت‌افزاری مکانیزم \lr{qcqa\_attention} روی مدل‌های پایه}
\label{tab:benchmark_qcqa_attention}
\begin{tabular}{lccccc}
\toprule
\textbf{مدل پایه} \lr{(Base Model)} & \textbf{دقت} \lr{(Accuracy)} & \textbf{شاخص} \lr{(F1-Score)} & \textbf{حافظه} \lr{(Peak VRAM)} & \textbf{توان} \lr{(Throughput)} & \textbf{وضعیت} \\
\midrule
\lr{DistilBERT} & \lr{63.76\%} & \lr{0.7299} & \lr{334.32 MB} & \lr{67,004.99 tok/s} & \lr{PASS} \\
\lr{RoBERTa} & \lr{51.03\%} & \lr{0.6748} & \lr{382.56 MB} & \lr{34,912.42 tok/s} & \lr{PASS\_LOW} \\
\lr{GPT-2} & \lr{50.92\%} & \lr{0.6738} & \lr{568.75 MB} & \lr{31,268.37 tok/s} & \lr{PASS\_LOW} \\
\lr{Qwen} & \lr{51.72\%} & \lr{---} & \lr{---} & \lr{---} & \lr{PASS} \\
\bottomrule
\end{tabular}
\end{table}

بر اساس نتایج تجربی، ماژول \lr{QCQA} بر روی مدل \lr{DistilBERT} دقت \lr{63.76\%} و شاخص F1 معادل \lr{0.7299} را به ثبت رسانده و اوج مصرف حافظه گرافیکی را در حد بسیار پایین \lr{334.32 MB} با توان عملیاتی خیره‌کننده \lr{67,004.99 tok/s} تثبیت کرده است.

\subsection{محدودیت‌ها و تحلیل شکست}
تحلیل تفصیلی کارنامه ماژول نمایانگر رفتارهای ساختاری زیر است:
\begin{itemize}
    \item \textbf{افت کارایی در مدل‌های \lr{RoBERTa} (\lr{51.03\%}) و \lr{GPT-2} (\lr{50.92\%}):} نتایج این دو مدل در آستانه \lr{PASS\_LOW} قرار دارد. علت بنیادین این پدیده، تداخل دینامیک وزن‌های مسیریاب با وزن‌های تثبیت‌شده پیش‌آموزش است. از آنجا که مقادیر اولیه مسیریاب به صورت تصادفی با وزن‌های نرمال کوچک مقداردهی می‌شوند، در فازهای نخستین آموزش بدون دوره گرم‌سازی طولانی، بازنمایی کلیدها دچار درهم‌آمیختگی و اغتشاش شده و مدل به جای بازیابی اطلاعات ظریف، به پیش‌بینی میانگین رده‌ها بسنده می‌نماید.
    \item \textbf{چالش مدل علّی \lr{Qwen} (\lr{51.72\%}):} در مدل‌های خودبازگشتی عمیق که از پیش بر پایه فشرده‌سازی صلب \lr{GQA} آموزش دیده‌اند، تغییر دینامیک فضاهای کلید/مقدار در طول دنباله مانع از همگرایی سریع گرادیان‌ها در لایه‌های فوقانی شبکه می‌گردد.
\end{itemize}
